% ══════════════════════════════════════════════════════════════════════════════
%  bench_integrators.tex
%  VSEPR-SIM Benchmark Suite — Layer 1
%  Velocity Verlet (NVE) vs Langevin-BAOAB (NVT) Integrator Comparison
%
%  Compile:  pdflatex bench_integrators.tex   (twice for TOC / refs)
%            or: latexmk -pdf bench_integrators.tex
%
%  Part of the FlowerOS HPC stack  (docs/bench_integrators.tex)
% ══════════════════════════════════════════════════════════════════════════════

\documentclass[11pt,a4paper]{article}

% ── packages ──────────────────────────────────────────────────────────────────
\usepackage[utf8]{inputenc}
\usepackage[T1]{fontenc}
\usepackage{lmodern}
\usepackage[margin=2.4cm]{geometry}
\usepackage{amsmath,amssymb,amsthm}
\usepackage{graphicx}
\usepackage{xcolor}
\usepackage{booktabs}
\usepackage{enumitem}
\usepackage{listings}
\usepackage{hyperref}
\usepackage{tikz}
\usetikzlibrary{arrows.meta,positioning,fit,calc,shapes.geometric}

% ── colours (FlowerOS pastel palette) ─────────────────────────────────────────
\definecolor{fosMint}{HTML}{87CEEB}
\definecolor{fosMagenta}{HTML}{D8A0D8}
\definecolor{fosGreen}{HTML}{87D687}
\definecolor{fosYellow}{HTML}{F0E68C}
\definecolor{fosGrey}{HTML}{A0A0A0}
\definecolor{fosBlue}{HTML}{6495ED}
\definecolor{fosOrange}{HTML}{FFA07A}
\definecolor{fosCode}{HTML}{F5F5F5}

\hypersetup{
  colorlinks=true,
  linkcolor=fosBlue,
  citecolor=fosBlue,
  urlcolor=fosMint,
  pdftitle={VSEPR Benchmark Layer 1 -- Integrators},
  pdfauthor={FlowerOS Project},
}

% ── listings style ────────────────────────────────────────────────────────────
\lstset{
  basicstyle=\ttfamily\small,
  backgroundcolor=\color{fosCode},
  frame=single,
  framerule=0.4pt,
  rulecolor=\color{fosGrey},
  keywordstyle=\color{fosBlue}\bfseries,
  commentstyle=\color{fosGrey}\itshape,
  stringstyle=\color{fosOrange},
  showstringspaces=false,
  breaklines=true,
  language=C++,
}

% ── theorem-like environments ─────────────────────────────────────────────────
\theoremstyle{definition}
\newtheorem{definition}{Definition}[section]
\newtheorem{contract}[definition]{Contract}
\theoremstyle{plain}
\newtheorem{proposition}[definition]{Proposition}

% ── macros ────────────────────────────────────────────────────────────────────
\newcommand{\VV}{\textsc{vv}}
\newcommand{\LG}{\textsc{lg}}
\newcommand{\NVE}{\textsc{nve}}
\newcommand{\NVT}{\textsc{nvt}}
\newcommand{\dt}{\Delta t}
\newcommand{\kb}{k_{\mathrm{B}}}
\newcommand{\half}{\tfrac{1}{2}}

% ── title ─────────────────────────────────────────────────────────────────────
\title{%
  \textbf{VSEPR-SIM Benchmark Suite} \\[4pt]
  \Large Layer 1: Integrator Performance Comparison \\[14pt]
  \normalsize Velocity Verlet (\NVE) \quad vs \quad Langevin-BAOAB (\NVT)%
}
\author{FlowerOS Project}
\date{v5.0.14 \quad \today}

% ══════════════════════════════════════════════════════════════════════════════
\begin{document}
\maketitle
\thispagestyle{empty}

\begin{abstract}
This document describes \textbf{Layer~1} of the VSEPR-SIM benchmark suite:
a wall-clock and accuracy comparison of two plugin integrators shipped in
\texttt{src/int/integrators.hpp}.
The \emph{Velocity Verlet} (\VV, microcanonical \NVE) integrator is compared
against the \emph{Langevin-BAOAB} (\LG, canonical \NVT) integrator over
log-spaced atom counts $N \in [32, 2048]$.
For each integrator and system size the benchmark records median time-per-step,
atom-step throughput, energy drift per atom (\VV\ only), and temperature
accuracy (\LG\ only).
All results are emitted as pipe-ready CSV to stdout and a human-readable
summary to stderr.
\end{abstract}

\tableofcontents
\newpage

% ──────────────────────────────────────────────────────────────────────────────
\section{Background and Motivation}
\label{sec:background}

Molecular dynamics (MD) simulation correctness and performance depend
critically on the choice of time integrator.  Two broad families exist:

\begin{description}[leftmargin=3em,labelwidth=2.5em]
  \item[\NVE] \emph{Microcanonical} ensemble: total energy $E = K + U$ is
	conserved.  The \emph{Velocity Verlet} algorithm is the canonical
	representative — it is symplectic (area-preserving in phase space),
	time-reversible, and second-order accurate in $\dt$.

  \item[\NVT] \emph{Canonical} ensemble: temperature $T$ is maintained by a
	thermostat.  The \emph{Langevin equation} adds a friction term and a
	Gaussian random force to Newton's equations, coupling the system to an
	implicit heat bath.  The BAOAB splitting scheme~\cite{baoab2012} is used
	here for its superior accuracy at moderate friction.
\end{description}

The FlowerOS VSEPR simulation stack exposes both integrators as plugin
objects behind a common step-interface, making a side-by-side comparison
straightforward.  This benchmark answers:

\begin{enumerate}
  \item How much per-step overhead does Langevin friction and RNG add relative
	to pure Velocity Verlet?
  \item At what system size does the $\mathcal{O}(N^2)$ force evaluation
	dominate and make the integrator overhead negligible?
  \item Does Velocity Verlet conserve energy (\emph{drift per atom}) to the
	expected level for the chosen time-step?
  \item Does Langevin maintain temperature to within a few percent after a
	short equilibration window?
\end{enumerate}

% ──────────────────────────────────────────────────────────────────────────────
\section{Integrator Algorithms}
\label{sec:algorithms}

\subsection{Velocity Verlet (NVE)}
\label{sec:vv}

Given positions $\mathbf{r}(t)$, velocities $\mathbf{v}(t)$, forces
$\mathbf{F}(t)$, and masses $m_i$, one Velocity Verlet step proceeds as:

\begin{align}
  \mathbf{v}(t + \tfrac{\dt}{2}) &= \mathbf{v}(t)
	+ \frac{\mathbf{F}(t)}{m_i}\,\frac{\dt}{2}
	\label{eq:vv-halfkick}\\
  \mathbf{r}(t + \dt) &= \mathbf{r}(t)
	+ \mathbf{v}(t + \tfrac{\dt}{2})\,\dt
	\label{eq:vv-drift}\\
  \mathbf{v}(t + \dt) &= \mathbf{v}(t + \tfrac{\dt}{2})
	+ \frac{\mathbf{F}(t+\dt)}{m_i}\,\frac{\dt}{2}
	\label{eq:vv-halfkick2}
\end{align}

where $\mathbf{F}(t+\dt)$ is evaluated from the updated positions.
The algorithm is symplectic and time-reversible; the global energy error
is $\mathcal{O}(\dt^2)$ per step, bounded $\mathcal{O}(\dt^2)$ in total
for any finite integration time (no secular drift for harmonic systems).

\subsection{Langevin-BAOAB (NVT)}
\label{sec:langevin}

The Langevin equation of motion is:
\begin{equation}
  m_i\,\ddot{\mathbf{r}}_i = \mathbf{F}_i
	- \gamma\,m_i\,\dot{\mathbf{r}}_i
	+ \sqrt{2\,\gamma\,m_i\,\kb T}\;\boldsymbol{\xi}(t)
  \label{eq:langevin}
\end{equation}
where $\gamma$ [ps$^{-1}$] is the friction coefficient, $T$ the target
temperature, and $\boldsymbol{\xi}$ Gaussian white noise with
$\langle\xi_\alpha(t)\xi_\beta(t')\rangle = \delta_{\alpha\beta}\,\delta(t-t')$.

The \textbf{BAOAB} splitting~\cite{baoab2012} factorises the propagator into:
\begin{enumerate}
  \item \textbf{B}: half-step velocity update from forces
  \item \textbf{A}: half-step coordinate drift
  \item \textbf{O}: full Ornstein--Uhlenbeck velocity relaxation
	$\mathbf{v} \leftarrow e^{-\gamma\dt}\mathbf{v}
	+ \sqrt{\kb T/m}\,(1-e^{-2\gamma\dt})^{1/2}\,\boldsymbol{\xi}$
  \item \textbf{A}: second half-step drift
  \item \textbf{B}: second half-step velocity update from new forces
\end{enumerate}

BAOAB is second-order accurate in $\dt$ for configurational observables
and exhibits smaller position autocorrelation error than the simpler
BBK~(Br{\"u}nger-Brooks-Karplus) scheme.

% ──────────────────────────────────────────────────────────────────────────────
\section{Benchmark Design}
\label{sec:design}

\subsection{Force Model}

Both integrators drive an identical $\mathcal{O}(N^2)$ Lennard-Jones
force field:
\begin{equation}
  U_{\mathrm{LJ}}(r) = 4\varepsilon\left[
	\left(\frac{\sigma}{r}\right)^{12}
	- \left(\frac{\sigma}{r}\right)^{6}\right]
  \label{eq:lj}
\end{equation}
with Ar-like parameters $\varepsilon = 0.1$\,kcal/mol,
$\sigma = 3.4$\,\AA, minimum distance clamp $r_{\min} = 0.5$\,\AA.
No cutoff is applied (all pairs evaluated), mirroring the VSEPR nonbonded
kernel baseline.

\subsection{System Generation}

Atoms are placed at uniform random positions in a $[-15, 15]^3$\,\AA\ box
with velocities drawn from $\mathcal{N}(0,\,0.5)$\,\AA/ps and masses
sampled uniformly from \{H\,1.008, C\,12.011, N\,14.007, O\,15.999\}\,u.

\subsection{Measurement Protocol}

\begin{enumerate}
  \item For each $N$ in a log-spaced sequence $[N_{\min}, N_{\max}]$:
  \begin{enumerate}
	\item Run $R$ independent repeats, each with a distinct RNG seed derived
	  from the base seed.
	\item Each repeat: build a fresh particle state, compute initial energy,
	  run exactly \texttt{n\_steps} MD steps (force~+ integration timed
	  together), compute final energy and temperature.
	\item Report the \emph{median} wall time across repeats.
  \end{enumerate}
\end{enumerate}

\subsection{Reported Metrics}

\begin{center}
\begin{tabular}{lll}
\toprule
Column & Formula & Applies to \\
\midrule
\texttt{median\_ms\_per\_step}
  & $\tilde{t}_{\mathrm{repeat}} / N_{\mathrm{steps}}$ & both \\
\texttt{atom\_steps\_per\_sec}
  & $N \cdot N_{\mathrm{steps}} / (\tilde{t}_{\mathrm{repeat}} \times 10^{-3})$
  & both \\
\texttt{energy\_drift\_per\_atom}
  & $|E_{\mathrm{final}} - E_{\mathrm{initial}}| / N$
  & \VV\ only \\
\texttt{temp\_error\_pct}
  & $|T_{\mathrm{actual}} - T_{\mathrm{target}}| / T_{\mathrm{target}} \times 100$
  & \LG\ only \\
\bottomrule
\end{tabular}
\end{center}

\subsection{CSV Schema}

\begin{lstlisting}[language={},basicstyle=\ttfamily\footnotesize]
n_atoms,integrator,median_ms_per_step,atom_steps_per_sec,
energy_drift_per_atom,temp_error_pct
\end{lstlisting}

Fields without a value for a given integrator are written as \texttt{n/a}.

% ──────────────────────────────────────────────────────────────────────────────
\section{Implementation Notes}
\label{sec:impl}

\subsection{Plugin Interface}

Both integrators share a two-call step pattern, allowing the force
evaluation to be injected between the first and second velocity half-kick:

\begin{lstlisting}
// Velocity Verlet
vv.step(coords, velocities, forces, dt, masses);   // half-kick + drift
eval_forces(state);                                // O(N^2) LJ
vv.step_second_half(velocities, forces, dt, masses); // second half-kick

// Langevin BAOAB
lang.step(coords, velocities, forces, dt, masses); // B-A-O-A
eval_forces(state);                                // O(N^2) LJ
lang.step_final(velocities, forces, dt, masses);  // final B
\end{lstlisting}

This design means the benchmark measures the \emph{true cost of one
complete MD step} including force evaluation.

\subsection{Energy and Temperature Calculations}

Kinetic energy:
\begin{equation}
  K = \sum_{i=1}^{N}\sum_{d \in \{x,y,z\}}
	  \half m_{i}\,v_{id}^{2}
\end{equation}

Instantaneous temperature from equipartition:
\begin{equation}
  T_{\mathrm{inst}}
	= \frac{2K}{(3N-3)\,\kb}
  \label{eq:temp}
\end{equation}
with $\kb = 1.987\times10^{-3}$\,kcal/mol/K.

\subsection{Build and Run}

The target is registered in \texttt{tests/CMakeLists.txt} as Group~11:
\begin{lstlisting}[language={},basicstyle=\ttfamily\small]
add_executable(bench_integrators bench_integrators.cpp)
target_link_libraries(bench_integrators vsepr_core vsepr_sim)
vsepr_add_test(NAME BenchIntegrators TARGET bench_integrators
			   LABELS benchmark integrator)
\end{lstlisting}

Invocation examples:
\begin{lstlisting}[language=bash,basicstyle=\ttfamily\small]
./bench_integrators                                  # stdout CSV + stderr summary
./bench_integrators --min 32 --max 2048 --repeats 5 # custom range
./bench_integrators --csv > integrators.csv          # pipe-only CSV
./bench_integrators --n_steps 500 --dt 0.0005        # tighter dt
\end{lstlisting}

% ──────────────────────────────────────────────────────────────────────────────
\section{Expected Results and Interpretation}
\label{sec:results}

\subsection{Throughput Scaling}

For small $N$, integrator overhead (RNG, friction update) is non-negligible
compared to the force evaluation.  As $N$ grows, the $\mathcal{O}(N^2)$ force
loop dominates and both integrators should converge to similar wall-time
per step.  The crossover typically occurs around $N \sim 200$–500 for a
scalar CPU implementation.

\subsection{Energy Drift (VV)}

For a time-step $\dt = 0.001$\,ps with LJ forces, energy drift per atom
should satisfy:
\begin{equation}
  \frac{\Delta E}{N} \lesssim \mathcal{O}\!\left(\varepsilon\,(\dt/\tau)^2\right)
\end{equation}
where $\tau \sim \sigma\,\sqrt{m/\varepsilon}$ is the LJ time unit.  Values
significantly above this indicate numerical instability (reduce $\dt$).

\subsection{Temperature Accuracy (Langevin)}

After $N_{\mathrm{steps}} \approx 100$ steps from a random configuration,
the system has not fully equilibrated.  Temperature errors of 10–40\% are
normal at short run times; the field represents convergence rate, not
steady-state error.  At long run times ($N_{\mathrm{steps}} \gtrsim 5000$)
errors below 5\% are expected for $\gamma = 1.0$\,ps$^{-1}$.

% ──────────────────────────────────────────────────────────────────────────────
\section{Defaults and CLI Reference}
\label{sec:cli}

\begin{center}
\begin{tabular}{llll}
\toprule
Flag & Type & Default & Description \\
\midrule
\texttt{--min}      & int    & 32    & Minimum atom count \\
\texttt{--max}      & int    & 2048  & Maximum atom count \\
\texttt{--steps}    & int    & 8     & Number of log-spaced size points \\
\texttt{--n\_steps} & int    & 100   & MD steps per timing window \\
\texttt{--repeats}  & int    & 3     & Independent repeats (median taken) \\
\texttt{--seed}     & uint32 & 42    & Base RNG seed \\
\texttt{--temp}     & double & 300.0 & Langevin target temperature (K) \\
\texttt{--friction} & double & 1.0   & Langevin friction $\gamma$ (ps$^{-1}$) \\
\texttt{--dt}       & double & 0.001 & Time-step (ps) \\
\texttt{--csv}      &        &       & Suppress stderr human summary \\
\texttt{--no-header}&        &       & Omit CSV header line \\
\bottomrule
\end{tabular}
\end{center}

% ──────────────────────────────────────────────────────────────────────────────
\section{Integration with the Benchmark Suite}
\label{sec:suite}

\begin{description}
  \item[Group 10 — GPU vs CPU (existing)] \texttt{bench\_gpu\_cpu}:
	compares scalar CPU and CUDA GPU for LJ energy evaluation.
  \item[\textbf{Group 11 — Integrators (this document)}]
	\texttt{bench\_integrators}: VV vs Langevin-BAOAB per-step cost and
	accuracy.
  \item[Group 12 — Spatial Kernels (companion document)]
	\texttt{bench\_spatial\_kernels}: direct $\mathcal{O}(N^2)$ vs
	simulation octree (see \texttt{docs/bench\_spatial\_kernels.tex}).
\end{description}

All three groups are reachable via CTest labels:
\begin{lstlisting}[language=bash,basicstyle=\ttfamily\small]
ctest -L benchmark          # all three groups
ctest -L integrator         # Group 11 only
ctest -L "benchmark;gpu"    # Group 10 only
\end{lstlisting}

% ──────────────────────────────────────────────────────────────────────────────
\begin{thebibliography}{9}

\bibitem{baoab2012}
B.~Leimkuhler and C.~Matthews,
\emph{Rational Construction of Stochastic Numerical Methods for Molecular
Sampling},
Appl.\ Math.\ Res.\ eXpress \textbf{2013}(1), 34--56 (2012).

\bibitem{swope1982}
W.~C.~Swope, H.~C.~Andersen, P.~H.~Berens, K.~R.~Wilson,
\emph{A computer simulation method for the calculation of equilibrium
constants for the association of integral membrane proteins},
J.\ Chem.\ Phys.\ \textbf{76}, 637 (1982).

\bibitem{frenkelsmit}
D.~Frenkel and B.~Smit,
\emph{Understanding Molecular Simulation: From Algorithms to Applications},
2nd~ed., Academic Press (2002).

\end{thebibliography}

\end{document}
